% Plan: development/outline.md, section 5, "§10".
% Run times: R/publication/integration/runtime-evidence-r1.json and
% R/publication/reproduction/commands.json (reproduction runs; catmix
% displayed bounds 794.65 or 2759.93, 1213.85, 2711.23, 2763.15 s; authors'
% catmix code 1049-4602 s; egbb and record/verify runs 194-1459 s);
% R/publication/reproduction/water-audit/logs/vbb2_*.time (63 fixed-target
% runs, 26,809 s wall, all "certified"), wTT_certify.log (regeneration);
% R/publication/reproduction/network/logs/ann.*.time (ANN replay 32:32 and
% 1:33:33 elapsed; verification 50:44 + 11:08 + 8:52 elapsed on two
% processes, user times 5,234 + 1,159 + 924 s = 2.03 h); development/reviews/
% (lnts, dtoc5 16.93 s, KAN 3,274 s); development/eg-audit/out/run_audit.log
% (44 jobs, 16,703 s summed, 2,215 s elapsed); the family dossiers and
% critiques (chain-catmix CC-08, waterno2 critique C1/C2, small dossier
% section 6). None of these times is in data/numbers.json; they support no
% speed claim. Per-result times are in the claim register (Appendix I /
% supplement). Regenerated waterno2 bounds: register N-39 and
% water-audit/logs/w18_certify.log, w24_certify.log. Relocation test:
% R/publication/reproduction/tools/smoke.py.
% Round-1 revision (adjudication G4-07): tiers are defined here only, by the
% recorded wall time per instance (sum over the parts of a split result).
% Per-instance sums checked 2026-10-04: water-audit/logs/vbb2_*.time gives
% waterno2_09 2,128 s, _12 2,944 s, _18 8,990 s, _24 12,747 s (63 runs,
% 26,809 s); development/eg-audit/out/run_audit.log (record + cert jobs) gives
% eg_int_s 465 s, eg_disc_s 1,456 s, eg_disc2_s 14,782 s (16,703 s in total).
% topopt: E-audit app:audit-existence (basis not stored; regenerated).
% Round-2 revision (adjudication G4-09), checked 2026-10-05: the regeneration
% details (waterno2_18/24 regenerated bounds 4790.820086 and 6576.150434, the
% clean-copy replay of 63 periods, powerflow0030p, the topopt basis) moved to
% app:repro-regen (move M6). KAN tiers by the per-instance rule, wall seconds
% of every code: r3 n4/n5/n9 first code (run_kan.py) 22.57/23.83/22.34
% (R/publication/integration/runtime-evidence-r1.json), unguarded second code
% 25.73/23.13/54.90 (development/dossiers/ann-kan-checks/logs/*.rigexp.log),
% guarded second code 24.74/22.71/58.86 (development/reviews/round1/kan-guard/
% report.md); r5 n3/n5/n8 first 1083.76/101.5/87.33, unguarded second
% 1221.20/760.55/1187.68, guarded 1175.59/744.42/1005.18. So r3 is Tier 1 and
% r5 Tier 2. eg_disc2_s: 14,782 s audited re-certification (Tier 3); 3.4 h is
% the CPU time of the original fast-mode search, 12,378 s (B7-eg, "CPU times").
% Round-3 revision (development/reviews/round3/), checked 2026-10-05: a result
% has the tier of the replay that certifies its displayed values (sol-referee
% 3); catmix displayed-bound replays 20-46 min (Tier 2), the other
% implementation (first code, catmix_bound.py) 1,049-4,602 s = 17-77 min
% (B4-chain-catmix certificate box; RUNS.md R06).
\section{Reproducibility and computational effort}\label{sec:repro}

All certificates, points, checkers and logs are in the archive that accompanies this paper (\archiveDOI).
The archive holds the stored models with their SHA-256 hashes (prefixes in \cref{tab:sem-hashes}), the certificate data (multipliers, split data, leaf boxes and coverage records, stage and pair bounds, PSD factors), the exact definitions of all exactly feasible points, the first implementations, the separately written implementations of \cref{sec:semantics-protocol} and the logs of every run cited here.
A number file stores, for every generated certified display (bounds, primal ends, gaps and derived margins), its exact value or certificate end, its source and the rounding direction; run times and fixed metadata come from the run records and the claim register (\cref{app:repro}).
The script that writes our tables checks the certified bound, primal, gap and margin displays that it generates against their exact values in rational arithmetic; certified displays in tables typed in the LaTeX sources and in the running text are checked separately (\cref{app:displays}).
A claim register maps each result to its checkers, recorded run time, tier, evidence level, status and trust-base terms, and the run records of the archive add the inputs and expected outputs (\cref{app:repro}).

\paragraph{Replay and regeneration.}
Replay (\cref{sec:semantics-protocol}) runs the checking code on stored inputs: multipliers, split data, partitions, targets and points.
At level~\evid{stored} it checks a stored finite certificate and never calls the numerical optimizer that proposed it; at level~\evid{rerun} it reruns, from the stored inputs, the deterministic search or computation whose results are not stored; where an LP solver takes part, it only proposes multipliers that are then checked rigorously.
A different LP build can change the tree, the run time and, in principle, whether the stored target is reached, but not the validity of a bound once it is certified.
Regeneration also reruns the numerical steps that produced the stored inputs (local solves, SDP and LP solves, solver-derived targets, time-limited searches) and may give a different valid certificate.
For \inst{waterno2_18}, \inst{waterno2_24} and \inst{powerflow0030p}, regeneration gives weaker valid bounds, and the reported values rest on the stored certificates.
The \inst{topopt-cantilever_60x40_50} existence certificate of the audit stores no basis, so it is regenerated rather than replayed; regeneration proved another exactly feasible point, which also satisfies its row of \cref{tab:audit-pairs}.
\Cref{app:repro-regen} gives the regenerated values.

\paragraph{Replay tiers.}
A replay belongs to Tier~1 if its recorded wall time per instance (for a result split into parts, the sum over its parts) is under 10 minutes, to Tier~2 if it is under one hour, and to Tier~3 otherwise.
A result has the tier of the replay that certifies its displayed values; where the other implementation's check takes longer, as for \inst{catmix}, its time is given separately.
\Cref{tab:repro-tiers} assigns the results; the claim register gives the recorded time of each.
The times come mostly from reproduction runs in a clean copy; some certificate boxes and the \inst{chain} run table (\cref{tab:chain-bnb}) quote other runs of the same or an equivalent check.
They show the order of magnitude to expect and support no speed comparison.
Several Tier~3 items split into periods, parts or chunks that can run in parallel.

\begin{table}[t]
\centering
\footnotesize
\caption{Replay tiers by recorded wall time per instance (Tier~1 under 10 minutes, Tier~2 under one hour, Tier~3 longer); one thread per process unless stated. Per-result times are in \cref{tab:repro-register}.}
\label{tab:repro-tiers}
\begin{tabular}{@{}>{\raggedright\arraybackslash}p{1.0cm}>{\raggedright\arraybackslash}p{13.4cm}@{}}
\toprule
tier & results and recorded times \\
\midrule
1 & \inst{lnts}, \inst{dtoc5}, \inst{optcdeg2}, \inst{camshape}, \inst{hvycrash}, \inst{pricing050}, \inst{etamac}, \inst{pindyck} (up to 322\,s), \inst{chain} (under 4\,min per instance and code), \inst{lukvle10} (about 5\,min), \inst{ex6_2_5}, \inst{ex6_2_7} (258\,s, 77\,s) and \inst{powerflow} bounds (\inst{powerflow0039p}, \inst{powerflow0039r} from the stored leaves); \inst{eg_int_s} (search and leaf re-certification under the accuracy auditor, 465\,s); KAN exact infeasibility; KAN bounds for $\RP$ of the three \texttt{r3} models (under 1\,min per model and code); exactly feasible points; audit existence proofs (\inst{topopt-cantilever_60x40_50} about 3\,min); the checks of \cref{sec:solvers}. \\
\addlinespace
2 & \inst{catmix} (replay of the displayed bounds, 20 to 46\,min each; the other implementation's check takes 17 to 77\,min); KAN bounds for $\RP$ of the three \texttt{r5} models (about 1.5 to 20\,min per model and code); \inst{eg_disc_s} (search and leaf re-certification under the accuracy auditor, 24\,min over two parts); \inst{waterno2_09}, \inst{waterno2_12} period bounds at the stored targets (35 and 49\,min). \\
\addlinespace
3 & \inst{waterno2_18}, \inst{waterno2_24} period bounds at the stored targets (2.5\,h and 3.5\,h); \inst{waterno2_06} pair bounds (49,315 bounds, about 60 CPU-hours on a loaded machine; the shortest path over them takes 11\,s); \inst{ann_cumene_tanh} (search replays 33\,min and 1.6\,h; verification 1.2\,h on two processes, 2.0 CPU-hours); \inst{eg_disc2_s} (re-certification of the stored leaves under the accuracy auditor, 4.1\,h over eight parts; its original search took about 3.4 CPU-hours). \\
\bottomrule
\end{tabular}
\end{table}

\paragraph{Disposable copies and environment.}
Every command runs in a disposable copy of the archive.
A relocation test reruns short checks in such a copy with the original tree hidden, records the files they open, and compares their output with the stored output (\cref{app:repro-setup}).
The checkers run under Python~3.13.11 with NumPy, SciPy (with HiGHS), mpmath and SymPy in the versions pinned in the archive (\cref{app:repro-setup}); PySCIPOpt (SCIP~10.0.2), CVXPY with Clarabel, and highspy propose multipliers, targets and candidate points but certify nothing.
The host is an Intel Xeon w5-2565X (18 cores) with 47\,GiB of memory under Ubuntu~24.04.4 in WSL2, shared with other jobs.
The package versions, CPU times and BLAS and libm dispatch of the original runs were not recorded consistently and are only partly recoverable.
The certificates do not rely on a particular build; \cref{tab:trust} states the trust base of each, and \cref{app:audit} that of the existence proofs of the audit.
